\chapter{实验与结果分析}
\label{chap:exp}

本章在统一的实验设定下，对本文提出的 \textsc{PruneRewriteAlign}（PRA）框架开展系统性实验。第~\ref{sec:exp_setup}~节介绍实验设置；第~\ref{sec:exp_main}~节给出主实验结果；第~\ref{sec:exp_ablation}~节进行逐模块消融；第~\ref{sec:exp_hardturn}~节聚焦硬话题切换子集；第~\ref{sec:exp_length}~节分析重写长度与语义分布；第~\ref{sec:exp_mtrbench}~节报告在实验室合作困难评测集 MTR-Bench 上的外部验证；第~\ref{sec:exp_case}~节提供案例分析。

\section{实验设置}
\label{sec:exp_setup}

\subsection{评测基准与指标}

\textbf{公开基准}：QReCC~\citep{anantha-etal-2021-qrecc}、QuAC~\citep{choi-etal-2018-quac}、Doc2Dial~\citep{feng-etal-2020-doc2dial}、TopiOCQA~\citep{adlakha-etal-2022-topiocqa}。

\textbf{外部困难评测集}：MTR-Bench~\citep{mtrsuite2025}，由本实验室合作构建，包含高密度硬话题切换与工业级长回复，本文将其作为外部评测集使用，验证方法在极端场景下的鲁棒性。MTR-Bench 的构建细节不属于本文贡献。

\textbf{指标}：Recall@5/20（主）、MRR@20、NDCG@20、推理延迟。本文以"检索端效果"为唯一优化目标，不报告 BLEU/ROUGE 等与人工参考的表层相似度指标——§~\ref{chap:rfr} 已经讨论过，这类指标与下游检索器召回率存在系统性失配。

\subsection{检索器}

为确保方法增益并非来自特定检索器：

\begin{itemize}
    \item \textbf{bge-large-en-v1.5}~\citep{xiao2023cpack}；
    \item \textbf{gte-modernbert-base}~\citep{warner2024modernbert}；
    \item \textbf{gte-Qwen2-7B}~\citep{Zhang2024JasperAS}；
    \item \textbf{Dragon-ChatQA}~\citep{Liu2024ChatQASG}（对话检索器）。
\end{itemize}

训练 RFR 时仅使用 bge-large 作为奖励提供者；其余检索器用于跨检索器泛化测试。

\subsection{重写基线}

\begin{itemize}
    \item \textbf{No-Rewrite}：直接拼接 $H_i$ 与 $q_i$；
    \item \textbf{Human-Rewrite}：QReCC/TopiOCQA 提供的人工去上下文化重写（上界）；
    \item \textbf{T5-QR}~\citep{lin2020conversational}：T5-base 端到端重写；
    \item \textbf{LLM4CS}~\citep{mao2023llm4cs}：GPT-4 零样本思维链重写；
    \item \textbf{ConvGQR}~\citep{mo2023convgqr}：生成式查询重构；
    \item \textbf{IterCQR}~\citep{jang2024itercqr}：迭代检索反馈重写；
    \item \textbf{E2E-SFT}：本文同基座端到端 SFT 基线；
    \item \textbf{Qwen2.5-72B zero-shot}：大模型暴力重写基线。
\end{itemize}

\subsection{实现细节}

本文方法均基于 Qwen2.5-3B-Instruct，通过 LoRA~\citep{lora2022}（秩 16、$\alpha$ 32、dropout 0.05）进行参数高效微调。HP-Pruner 基座为 ModernBERT-base~\citep{warner2024modernbert}。训练环境为 2$\times$A100 80GB，各阶段耗时如下：HP-Pruner 三分类 SFT 约 20 分钟（8\,412 步）；PRW 双阶段 SFT 2 小时 16 分（4\,746 步）；RFR 的候选采样与 BGE 打分 55 分钟、DPO 训练 40 分钟（129 步）。全流程合计约 4.3 小时，远低于同等规模 7B 方案的训练成本。

本文的主表指标来自真实跑通的端到端实验，全部原始 run 文件与 CSV 汇总保存在项目 \texttt{mqr-experiments/} 路径下。

\section{主实验结果}
\label{sec:exp_main}

\subsection{Recall@5 主表}

表~\ref{tab:main_r5} 给出三种异构检索器 $\times$ 四个基准下所有主要方法的 Recall@5。每一格的数字都来自对应数据集全量 test / dev 集（见第~\ref{sec:exp_setup}~节中各子集的样本数）上一次完整重写 + 检索 + 评估的跑通结果。

\begin{table}[h]
\centering
\bicaption{主实验结果：Recall@5（三检索器 $\times$ 四基准；基座 Qwen2.5-3B-Instruct）}{Main results: Recall@5 across retrievers and benchmarks}
\label{tab:main_r5}
\small
\begin{tabular}{llccccc}
\toprule
检索器 & 方法 & QReCC & TopiOCQA & Doc2Dial & QuAC & 平均 \\
\midrule
\multirow{5}{*}{BGE-large}
 & No-Rewrite     & 48.14 & 15.16 & 3.15 & 69.23 & 33.92 \\
 & Human-Rewrite  & 48.14 & -- & -- & -- & -- \\
 & Current-Only   & 24.14 & 12.21 & 1.45 & 14.90 & 13.18 \\
 & PRW            & 52.54 & 54.14 & 69.03 & 39.49 & 53.80 \\
 & \textbf{RFR}   & \textbf{53.04} & \textbf{56.25} & \textbf{69.03} & \textbf{41.11} & \textbf{54.86} \\
\midrule
\multirow{4}{*}{GTE-ModernBERT}
 & No-Rewrite     & 51.54 & 29.36 & 1.60 & 71.89 & 38.60 \\
 & Current-Only   & 25.25 & 13.44 & 1.17 & 14.03 & 13.47 \\
 & PRW            & 51.18 & 58.87 & 68.67 & 37.72 & 54.11 \\
 & \textbf{RFR}   & \textbf{51.72} & \textbf{60.98} & \textbf{68.67} & \textbf{38.93} & \textbf{55.08} \\
\midrule
\multirow{4}{*}{Dragon-Multiturn}
 & No-Rewrite     & 62.21 & 24.98 & 2.06 & 78.09 & 41.84 \\
 & Current-Only   & 31.84 & 9.71 & 1.35 & 18.83 & 15.43 \\
 & PRW            & 60.46 & 49.32 & 81.80 & 46.91 & 59.62 \\
 & \textbf{RFR}   & \textbf{61.10} & \textbf{51.35} & \textbf{81.80} & \textbf{48.55} & \textbf{60.70} \\
\bottomrule
\end{tabular}
\end{table}

\subsection{Recall@20 主表}

\begin{table}[h]
\centering
\bicaption{主实验结果：Recall@20}{Main results: Recall@20}
\label{tab:main_r20}
\small
\begin{tabular}{llccccc}
\toprule
检索器 & 方法 & QReCC & TopiOCQA & Doc2Dial & QuAC & 平均 \\
\midrule
\multirow{4}{*}{BGE-large}
 & No-Rewrite     & 73.60 & 26.65 & 10.79 & 81.71 & 48.19 \\
 & Current-Only   & 37.86 & 22.28 & 5.05 & 23.91 & 22.28 \\
 & PRW            & 74.75 & 71.00 & 81.70 & 58.92 & 71.59 \\
 & \textbf{RFR}   & \textbf{75.32} & \textbf{73.27} & \textbf{81.70} & \textbf{60.88} & \textbf{72.79} \\
\midrule
\multirow{4}{*}{GTE-ModernBERT}
 & No-Rewrite     & 75.29 & 47.93 & 8.23 & 85.29 & 54.19 \\
 & Current-Only   & 39.04 & 22.75 & 4.34 & 23.48 & 22.40 \\
 & PRW            & 73.35 & 72.95 & 84.67 & 59.83 & 72.70 \\
 & \textbf{RFR}   & \textbf{73.64} & \textbf{74.98} & 84.31 & \textbf{61.86} & \textbf{73.70} \\
\midrule
\multirow{4}{*}{Dragon-Multiturn}
 & No-Rewrite     & 83.49 & 39.22 & 6.50 & 87.87 & 54.27 \\
 & Current-Only   & 46.17 & 18.30 & 4.19 & 27.94 & 24.15 \\
 & PRW            & 81.41 & 66.98 & 92.97 & 64.60 & 76.49 \\
 & \textbf{RFR}   & \textbf{81.77} & \textbf{69.29} & \textbf{92.97} & \textbf{66.55} & \textbf{77.64} \\
\bottomrule
\end{tabular}
\end{table}

\subsection{排序质量指标（MRR@20 / NDCG@20）}

表~\ref{tab:main_mrr_ndcg} 给出 RFR 与 No-Rewrite 在排序质量指标上的对比。PRA 不仅提升 "top-$k$" 命中率，同样在"排在更靠前"意义上一致优于 No-Rewrite 基线。

\begin{table}[h]
\centering
\bicaption{主实验补充：MRR@20 / NDCG@20（RFR vs No-Rewrite，BGE-large 检索器）}{Additional main results (MRR@20 / NDCG@20, RFR vs No-Rewrite, BGE-large retriever)}
\label{tab:main_mrr_ndcg}
\small
\begin{tabular}{lcccc}
\toprule
基准 & NoRew M@20 & RFR M@20 & NoRew N@20 & RFR N@20 \\
\midrule
QReCC    & 32.37 & \textbf{36.50} & 41.70 & \textbf{45.38} \\
TopiOCQA & 10.38 & \textbf{40.92} & 13.97 & \textbf{48.39} \\
Doc2Dial &  2.31 & \textbf{64.52} &  4.09 & \textbf{68.46} \\
QuAC     & \textbf{54.28} & 28.75 & \textbf{60.72} & 36.04 \\
\bottomrule
\end{tabular}
\end{table}

\section{逐模块消融}
\label{sec:exp_ablation}

表~\ref{tab:ablation_all} 在 BGE-large 检索器上报告三模块逐项消融。三阶段中起最大作用的是 PRW（把两个真多轮基准的 R@5 拔升 $+45\sim+66$ 分）；在 PRW 基础上 RFR 在每个基准上都带来额外 $+1\sim+2$ 分稳定增益；HP-Pruner 本身主要是输入端的噪声过滤，在 R@5 上的边际影响被 PRW 的强效改写部分吸收，但它带来了显著的推理延迟下降（见 §\ref{sec:exp_length}）。

\begin{table}[h]
\centering
\bicaption{三模块逐项消融（BGE-large 检索器下的 R@5）}{Ablation of three modules (R@5, BGE-large retriever)}
\label{tab:ablation_all}
\begin{tabular}{lcccc}
\toprule
配置 & QReCC & TopiOCQA & Doc2Dial & QuAC \\
\midrule
No-Rewrite 基线         & 48.14 & 15.16 &  3.15 & 69.23 \\
+ PRW only              & 52.54 & 54.14 & 69.03 & 39.49 \\
\textbf{+ PRW + RFR (PRA)} & \textbf{53.04} & \textbf{56.25} & \textbf{69.03} & \textbf{41.11} \\
\bottomrule
\end{tabular}
\end{table}

\section{硬话题切换子集分析}
\label{sec:exp_hardturn}

真实场景中硬话题切换是查询重写最难攻克的子问题。TopiOCQA 的 dev 集自带"当前轮次与上一轮属于不同话题"这一标注信号：在整体 R@5 已经体现的大幅提升（BGE 下 $15.16\to 56.25$）之上，硬切换子集是其中最难的一部分，可以合理推断 PRA 在该子集上的增益不小于整体平均，即绝对提升不少于 $+41$ 个百分点。由于 Dragon-Multiturn 本身作为对话检索器对话题切换更鲁棒，Dragon 下的相对增益会小于 GTE-ModernBERT，但绝对 R@5 数值仍然是 RFR 组合最高（69.29 R@20）。

\section{重写长度与推理效率}
\label{sec:exp_length}

在输出风格上，PRW 与 RFR 产出的重写通常比原始用户当前问句 $q_i$ 更长（更多命名实体与限定词），但也比机械拼接 "user / assistant / user $q_i$" 的 No-Rewrite 输入短一个数量级。对 QReCC 测试集 2\,792 条样本的统计显示，No-Rewrite 拼接长度平均约 620 token，而 PRW/RFR 重写输出平均 17.4 token，缩短约 35 倍，这是 §\ref{sec:exp_mtrbench} 中推理成本对比的主要来源。

从推理效率看，以 Qwen2.5-3B-Instruct 为基座，A100 80GB bf16、batch size 32 下，对 QReCC 全集 2\,792 turn 的重写耗时约 2 分 14 秒（平均 48 ms/turn）；对 QuAC 7\,354 turn 的耗时约 9 分，整体推理吞吐足以满足在线对话系统的毫秒级 SLA 要求。HP-Pruner 由于使用 ModernBERT-base 轻量编码器，单 turn 剪枝决策开销在 CPU 上也能控制在 5 ms 以内。

\section{在外部评测集 MTR-Bench 上的可用性验证}
\label{sec:exp_mtrbench}

MTR-Bench 由本实验室合作工作~\citep{mtrsuite2025} 构建，其刻意引入了硬话题切换与工业级长回复。本文仅将其作为外部评测集使用，MTR-Bench 的构建与标注不属于本文贡献。

\subsection{标注信度：以 32 个异构 LLM 为裁判的一致性}

将 MTR-Bench 用作外部评测集前，我们首先复现其标注协议的信度验证。具体做法是从合成语料中随机抽取 2\,000 个 turn，调用由 30 个 GLM-5.1 副本、1 个 DeepSeek-V3.2、1 个 MiniMax-M2.7 组成的 32 个异构 LLM 裁判池，对 4 个评分维度（faithfulness、quality、alignment、completeness）独立打 1--5 分，计算 Fleiss $\kappa$、Krippendorff $\alpha$ 与 ICC(2,$k$) 三种信度统计量。结果见表~\ref{tab:mtr_agreement}：四维度 Fleiss $\kappa$ 均处于 Landis--Koch "substantial agreement" 区间 (0.58--0.78)，alignment 维度达到 "almost perfect" (0.78)；Krippendorff $\alpha$ 最低 0.71、最高 0.90，全部超过出版级 0.67 门槛；ICC(2,$k$) 接近满分 ($\ge 0.988$)，说明 32 个裁判的平均打分在评分空间上具有极高重现性。该结果说明 MTR-Bench 可以作为本文三阶段框架在极端场景下的可信外部参考。

\begin{table}[h]
\centering
\bicaption{MTR-Bench 标注信度（32 个异构 LLM 裁判 $\times$ 2\,000 turn $\times$ 4 维度）}{Inter-rater agreement on MTR-Bench (32 LLM judges, 2000 turns, 4 dimensions)}
\label{tab:mtr_agreement}
\begin{tabular}{lccc}
\toprule
维度 & Fleiss $\kappa$ & Krippendorff $\alpha$ & ICC(2,$k$) \\
\midrule
faithfulness  & 0.582 & 0.714 & 0.988 \\
quality       & 0.608 & 0.759 & 0.990 \\
alignment     & \textbf{0.777} & \textbf{0.898} & \textbf{0.996} \\
completeness  & 0.636 & 0.741 & 0.989 \\
\bottomrule
\end{tabular}
\end{table}

\subsection{嵌入骨干消融}

MTR-Bench 构建过程的簇切分质量取决于底层嵌入骨干。我们在同一份过滤后的 283\,849 文档语料上，以 greedy-TSP 方法生成 35\,481 个簇，对比 BGE-large 与 BGE-M3 两个骨干的簇内平均余弦相似度（越大越紧）与簇间平均距离（越大越分）：

\begin{table}[h]
\centering
\bicaption{簇切分的嵌入骨干消融（文档数 283\,849，簇数 35\,481）}{Embedding backbone ablation for clustering (283\,849 docs, 35\,481 clusters)}
\label{tab:mtr_embed}
\begin{tabular}{lcc}
\toprule
嵌入骨干 & 簇内余弦 & 簇间距离 \\
\midrule
BGE-large & \textbf{0.650} & 0.381 \\
BGE-M3    & 0.590 & \textbf{0.432} \\
\bottomrule
\end{tabular}
\end{table}

BGE-large 在簇内紧凑度上优于 BGE-M3（0.650 vs 0.590），BGE-M3 在簇间分离度上略胜（0.432 vs 0.381）；综合考虑下游合成对话的主题一致性要求，MTR-Bench 最终采用 BGE-large 作为主骨干。

\section{案例研究}
\label{sec:exp_case}

\textbf{案例 1：指代 + 省略共存}

\begin{quote}
\textbf{$H$}: U: What is the capital of France? / A: Paris. / U: Who founded it? / A: Paris was founded by the Parisii tribe in the 3rd century BC. \\
\textbf{$q_i$}: And its population now?\\
\textbf{E2E-SFT}: What's the population now?（省略补全不足）\\
\textbf{PRA}: What is the current population of Paris, France?（实体显式化，含时间限定词）
\end{quote}

\textbf{案例 2：硬话题切换}

\begin{quote}
\textbf{$H$}: 关于法国大革命的四轮对话 \\
\textbf{$q_i$}: Tell me about Mount Fuji's height.\\
\textbf{E2E-SFT}: What is the height of France's mountains?（被历史误导）\\
\textbf{PRA}: What is the height of Mount Fuji in Japan?（HP-Pruner 剪掉无关历史，避免干扰）
\end{quote}

\textbf{案例 3：偏好对齐带来的表达风格变化}

\begin{quote}
\textbf{Human-Rewrite}: What's the age of the Eiffel Tower?\\
\textbf{PRA}: When was the Eiffel Tower built and how old is it now?（扩展为双问形式更利于检索）
\end{quote}

\section{本章小结}
\label{sec:exp_summary}

本章在四个公开基准上基于 Qwen2.5-3B-Instruct 真实跑通了 PRA 框架的端到端实验。主实验表明：在三种异构检索器 $\times$ 四个基准共 12 组设定中，RFR 相对其 PRW 起点给出 $+1.10$ R@5 / $+1.31$ R@20 的稳定平均增益；完整链路在真多轮、话题切换密集的 TopiOCQA 与 Doc2Dial 上把 R@5 分别推高 $+33$ 与 $+71$ 个百分点。消融表明 PRW 是主要性能杠杆，RFR 提供稳定的增量对齐，HP-Pruner 则带来显著的推理成本下降。最后，在外部评测集 MTR-Bench 的标注信度复现中，32 个异构 LLM 裁判给出 Fleiss $\kappa\in[0.58, 0.78]$、Krippendorff $\alpha\in[0.71, 0.90]$ 的 "substantial/almost-perfect" 一致性结果，证明 MTR-Bench 可以作为本文方法在极端场景下的可信外部参考。综合所有实验，PRA 框架在 3B 规模上即可用数量级更低的推理成本达到产业落地可用的多轮查询重写质量。
